\chapter*{Abstract}\label{abstract}
Robotic perception often requires identifying not only what object is present, but which viewpoint best satisfies a task-specific natural language instruction. Different views of the same object reveal different visual features such as handles, printed designs, or affordance regions. Under a limited observation budget, a system must therefore select a small number of informative views rather than inspect every candidate.

This thesis investigates language-guided active view selection for multi-view object perception. The system starts from a random initial view and sequentially selects additional views under budgets \(k \in \{1,2,3,4,5\}\). After the budget is exhausted, an observed-view reranker returns a final prediction. The work evaluates a CLIP relevance-plus-novelty active policy, a learned MLP reranker trained on frozen CLIP features, and a trained active view selector that uses the same CLIP+MLP architecture. These methods are compared with random selection and exhaustive CLIP ranking.

Two datasets are used. The first is a CO3D-based multi-view annotation set covering five object categories with 30 instances each and five viewpoints per instance. Ground-truth views are instruction-dependent and view order is shuffled to remove first-view bias. The second is derived from the UMD+GT RGB-D affordance dataset. Affordance masks are converted into language-guided view-selection samples, where ground-truth views are those whose affordance area is at least 80\% of the largest area among the five candidates.

On the shuffled CO3D evaluation set, learned reranking is the main source of improvement: at budget~3, random selection with learned reranking reaches 0.484 accuracy, while the full trained-selection plus learned-reranking pipeline reaches 0.520, compared with 0.310 for the random-plus-CLIP baseline. On UMD+GT, evaluated over 20 random-initial trials, the full pipeline is clearly strongest, reaching \(0.669\pm0.011\) already at budget~2 and \(0.680\pm0.008\) at budget~3, against about 0.46 for random-plus-CLIP. Ablation studies show that the learned reranker contributes most of the gain, while the trained selector is especially useful when the dataset contains smoothly varying viewpoint cues, as in UMD affordance sequences.

Overall, this thesis shows that language-guided scoring can support active view selection, that a lightweight MLP on frozen CLIP features substantially improves final view ranking, and that a trained selector can further improve low-budget performance when viewpoint variation is informative. Remaining challenges include overlapping selector/reranker training objectives, optimistic access to unseen-view features during selection, and limited sample counts for some affordance categories.
